% ============================================
% LITERATURE REVIEW
% ============================================

\begin{frame}{Literature Review}
    \begin{table}
        \centering
        \small % अक्षर ठिक्कको साइजमा राख्न
        \caption{Summary of existing skyline geolocation approaches}
        \vspace{0.5em}
        \resizebox{\textwidth}{!}{%
            \begin{tabular}{lllc}
                \toprule
                \textbf{Authors} & \textbf{Year} & \textbf{Terrain} & \textbf{Matching Method} \\
                \midrule
                Baboud et al. & 2011 & Alps & Edge cross-correlation \\
                Baatz et al.  & 2012 & Alps (Switzerland) & Bag-of-curves \\
                Saurer et al. & 2016 & Alps & Contours + DTW \\
                Pan et al.    & 2020 & China (hilly) & Euclidean + heatmap \\
                Tomešek et al.& 2022 & Alps & Synthetic rendering \\
                \midrule
                \textbf{This Work} & \textbf{2026} & \textbf{Khumbu, Nepal} & \textbf{FFT + DTW} \\
                \bottomrule
            \end{tabular}
        }
    \end{table}
    
    \vspace{1em}
    \begin{itemize}
        \item \small \textbf{Research Gap:} Existing systems focus on European/Chinese terrain; none address the extreme elevation and GPS failure challenges in the \textbf{Himalayas}.
    \end{itemize}
    
\end{frame}